% !TeX root = ../../../../../main.tex
% chapters/sections/chapter5/subsections/beta_shock/summary.tex

\subsection{まとめ}
\label{sec:chapter5_beta_shock_summary}
本節の締めくくりとして、需要ショック（\(\beta\)ショック）における各金融政策のパフォーマンスの違いをもたらした要因を総括し、あわせて名目総消費水準目標（NCLT）が全政策中で最高評価（厚生\(W_s^H\)の落ち込みが最小）を獲得した理由について考察する。

家計が貯蓄志向を高める\(\beta\)ショックにおいては、現在需要の急減に伴う深刻な不況とデフレ圧力が同時に発生し、名目利子率が速やかにゼロ金利制約（ZLB）に直面する。この厳しい制約下においては、中央銀行が「いかに早く実体経済の崩壊を食い止めるか（初動の速さ）」と、「回復期においていかに適切な物価のアンカーを保つか（名目アンカー）」に対する姿勢の違いが、各政策の明暗を大きく分けることとなった。

\begin{enumerate}
    \item
    \label{itm:chapter5_beta_shock_summary_real_target}
        \textbf{初動に優れるが名目アンカーを欠く実質目標群の限界}：
        総消費水準目標（CLT）や生産水準目標（OLT）は、ジャンプ変数である実質変数を直接目標としているため、ショック発生と同時に名目利子率をゼロまで引き下げ、前半の実体経済の落ち込みを全政策の中で最も小さく防ぐことに成功した。しかし、これらの政策は名目アンカーを持たないため、中盤以降の物価の果てしない下落ドリフトに歯止めをかけることができず、インフレ率の極端な変動とそれに伴う価格分散（\(\Delta_{t-1}^H\)）の際限のない爆発的な拡大を招いてしまい、最高評価には至らなかった。
    \item
    \label{itm:chapter5_beta_shock_summary_price_target}
        \textbf{物価安定に優れるが初動が遅れる物価目標群の敗因}：
        対照的に、生産者物価インフレ目標（IT-PPI）や生産者物価水準目標（PPLT）は、物価と価格分散を完全に安定させたものの、目標変数が遅行指標である物価のみに依存している。そのため、ショック直後に実体経済を支えるための素早い利下げ（初動）に遅れをとり、初期の消費と生産の深刻な大暴落を防ぐことができなかった。さらに、消費者物価を目標とする政策群（CPLT, IT-CPI, TR-CPI）に至っては、下落した物価を強引に引き戻そうとして中盤以降に巨大な波打ち（消費のオーバーシュート）を引き起こし、マクロ経済を極度に不安定化させたため最低評価に沈む結果となった。
\end{enumerate}

\paragraph{名目総消費水準目標（NCLT）が最高評価を得た理由}
\label{sec:chapter5_beta_shock_summary_nclt_advantage}
上記のようなゼロ金利制約下における厳しい環境において、NCLTが最も高い厚生評価を獲得したのは、各政策群が持つ長所を単一のルール内で極めて高次元にハイブリッドさせているためである。

NCLTは「名目総消費（\(p_t^{H\to W}c_t^{H\to W}\)）」を目標としている。目標の中に実質消費（ジャンプ変数）を含んでいるため、ショック発生時にはCLTと同様に素早く利子率をゼロに張り付け、前半の実体経済の落ち込みを最小限に食い止めることができる（\textbf{前半の強さ}）。
さらに、中盤以降の回復期においては、CPLTやIT-CPIが持つ「後半の巨大な消費増加」という長所を適度な形で取り入れており、CLTが速やかに定常状態に戻ってしまうのに対し、NCLTはより長く高い水準で消費を維持し続けることで厚生を押し上げている（\textbf{後半の強さ}）。
それと同時に、目標の中に物価水準を含んでいるため、経済が回復し物価が変動し始めた際には名目アンカーとして適切にブレーキがかかる。これにより、CLTが陥った「物価ドリフトと価格分散の爆発」を未然に防ぎ、価格分散を極めて低い水準に抑え込みながら経済を定常状態へと軟着陸させているのである。

このように、NCLTは「消費という実質変数による鋭い初動」、「長期的な消費の高水準維持」、および「名目枠の縛りによる価格分散コストの抑制」という3つの要素を見事に同時に達成している。供給ショック（\(a\)ショック）における分析結果と併せると、NCLTはいかなる経済ショックに対しても、極端な不況を回避しつつ物価の安定を担保できる、極めて優れた頑健性（ロバストネス）を持つ金融政策ルールであると結論付けることができる。